% ----------------------------------------------------------------------
%  3.5.2. Зависимость электропроводности металлов от температуры;
%  3.5.3. Эффект Холла
% ----------------------------------------------------------------------
\subsection{Зависимость электропроводности металлов от температуры}\label{ssec:rho-T}

Сопротивление металлов, как правило, растёт с температурой. Электроны
рассеиваются двумя независимыми способами:
\begin{itemize}
\item на \emph{статических} центрах~— дефектах решётки, атомах примесей;
\item на \emph{тепловых колебаниях} решётки~— фононах.
\end{itemize}
Частоты независимых столкновений складываются,
$1/\tau=1/\tau_\text{деф}+1/\tau_\text{фон}$, откуда \term{правило Маттиссена}:
\begin{keyeq}\label{eq:matthiessen}
\rho=\rho_\text{ост}+\rho_\text{ид}(T).
\end{keyeq}
\term{Остаточное сопротивление} $\rho_\text{ост}$ не зависит от
температуры и тем меньше, чем чище и совершеннее образец;
\term{идеальное} $\rho_\text{ид}(T)$ создаётся фононами и сильно зависит
от температуры. Отношение сопротивлений при комнатной и гелиевой
температурах $R_{293}/R_{4{,}2}$ служит мерой чистоты и совершенства
металла; у совершенных высокочистых монокристаллов оно достигает
$10^5$.

\paragraph{Высокие температуры ($T>\theta_D$)}
Здесь главное~— рассеяние на фононах ($\theta_D$~— температура Дебая,
п.~\ref{sec:thermal}).
\begin{shagi}
\item Колеблющийся атом можно представить диском радиусом в амплитуду
колебаний $a$; сечение рассеяния $S=\pi a^2$, а длина пробега $l\propto1/a^2$.
\item Энергия колебаний $\propto a^2$ и пропорциональна температуре,
поэтому $a^2\propto T$ и $l\propto1/T$. (То же следует из того, что число
фононов при высокой температуре $\propto T$.)
\item Фононы при высоких температурах имеют большие импульсы, и одного
столкновения хватает, чтобы электрон потерял направленную скорость
($\nu\approx1$).
\end{shagi}
\noindent Отсюда $\sigma\propto T^{-1}$, т.\,е. сопротивление растёт
\emph{линейно}:
\begin{keyeq}\label{eq:rho-alpha}
\rho_t=\rho_0\bigl[1+\alpha(t-t_0)\bigr],
\end{keyeq}
где $\alpha$~— \term{температурный коэффициент сопротивления}. Для
непереходных металлов $\alpha\approx0{,}004$~К$^{-1}$, у переходных и
особенно ферромагнитных~— больше (до 0,01~К$^{-1}$), см.
табл.~\ref{tab:rho-metals}.

\begin{table}[ht]
\centering
\caption{Удельное сопротивление и температурный коэффициент при 0\,°C,
температура Дебая}
\label{tab:rho-metals}
\small
\begin{tabular}{@{}lccccccc@{}}
\toprule
 & Al & Be & W & Fe & Cu & Mo & Ti \\
\midrule
$\rho$, мкОм$\cdot$см & 2,50 & 3,2 & 4,89 & 8,6 & 1,55 & 5,03 & 42 \\
$\alpha$, $10^{-5}$~К$^{-1}$ & 460 & 900 & 510 & 651 & 433 & 473 & 546 \\
$\theta_D$, К & 433 & 1481 & 383 & 477 & 347 & 423 & 420 \\
\bottomrule
\end{tabular}
\end{table}

\paragraph{Низкие температуры: примеси}
При $T\ll\theta_D$ фононов мало, и в не очень чистых металлах главным
становится рассеяние на примесях и дефектах. Длина пробега здесь
определяется концентрацией примесей и от температуры не зависит; а так
как в рассеянии участвуют лишь фермиевские электроны со скоростью
$v_\text{F}$, тоже не зависящей от температуры, то и проводимость
постоянна. Отсюда~— температурно-независимое остаточное сопротивление в
правиле Маттиссена.

\paragraph{Низкие температуры: фононы, закон $T^5$}
В очень чистых металлах и при низких температурах главным остаётся
рассеяние на фононах, но законы иные.
\begin{shagi}
\item Число фононов при $T\ll\theta_D$ пропорционально $T^3$ (теория
Дебая, п.~\ref{sec:thermal}), поэтому длина пробега
$l\propto1/n_\text{ф}\propto T^{-3}$\ispr{В оригинале (с.~500) записано
$l\sim1/n_\text{ф}\sim T^3$; правильно $l\propto T^{-3}$.}.
\item Импульс фонона мал по сравнению с импульсом электрона, и одно
столкновение лишь слегка поворачивает электрон на угол
$\alpha\approx p_\text{ф}/p_\text{эл}$. Направленный импульс уменьшается
на $\Delta p\approx p_\text{эл}(1-\cos\alpha)\approx p_\text{эл}\alpha^2/2$.
При низких температурах $(p_\text{ф}/p_\text{эл})^2\sim T/\theta_D$.
\item Чтобы погасить импульс, нужно $\nu\approx p_\text{эл}/\Delta p\propto(\theta_D/T)^2$
столкновений, т.\,е. $\nu\propto T^{-2}$.
\item Итого по \eqref{eq:sigma-quantum}: $\sigma\propto\nu l\propto T^{-2}T^{-3}=T^{-5}$,
т.\,е. $\rho\propto T^5$.
\end{shagi}
Точная теория даёт \term{формулу Блоха--Грюнайзена}
\begin{equation}\label{eq:bloch-gruneisen}
\rho_\text{фон}=\frac{K}{A\theta_D}\left(\frac T{\theta_D}\right)^5
\int_0^{\theta_D/T}\frac{x^5\,\d x}{(e^x-1)(1-e^{-x})}
\end{equation}
($K$~— постоянная, $A$~— атомная масса)\ispr{В оригинале (с.~501)
подынтегральное выражение записано как $x^4/(e^x-1)$. С ним предел
интеграла при $T\to0$ равен $\approx24{,}9$, а не приведённому там же
числу $124{,}4$. Число $124{,}4=5!\,\zeta(5)$ соответствует
стандартному виду формулы Блоха--Грюнайзена с $x^5/[(e^x-1)(1-e^{-x})]$;
высокотемпературный предел при этом прежний.} с двумя пределами
(рис.~\ref{fig:bloch-gr}):
\[
T\ge\theta_D:\ \rho_\text{фон}\approx\frac{KT}{4A\theta_D^2};\qquad
T\ll\theta_D:\ \rho_\text{фон}\approx124{,}4\,\frac{KT^5}{A\theta_D^6}.
\]
Для металлов с почти сферической поверхностью Ферми (щелочных) закон
$T^5$ выполняется уже при $T/\theta_D<0{,}1$, для металлов со сложной
поверхностью Ферми~— при гораздо более низких температурах. У бериллия
$\theta_D$ очень высока, и нелинейный участок тянется в широком
интервале температур.

\begin{figure}[ht]
\centering
\begin{tikzpicture}
\begin{axis}[width=0.66\textwidth, height=5.4cm, xmin=0, xmax=1.2, ymin=0, ymax=1.35,
  xlabel={$T/\theta_D$}, ylabel={$\rho/\rho_\text{фон}(\theta_D)$}, axis lines=left,
  legend style={font=\scriptsize, at={(0.03,0.97)}, anchor=north west, draw=none, fill=none}]
  \addplot[fmGray, dashed] table[x=t, y=lin] {figures/el-bg.dat};
  \addlegendentry{асимптота $\rho\propto T$}
  \addplot[fmBlue, thick] table[x=t, y=r] {figures/el-bg.dat};
  \addlegendentry{чистый металл}
  \addplot[fmRed, thick] table[x=t, y=r2] {figures/el-bg.dat};
  \addlegendentry{с примесями ($+\rho_\text{ост}$)}
  \node[font=\scriptsize, anchor=west] at (axis cs:0.08,0.18) {$\rho\propto T^5$};
  \draw[-{Stealth[length=1.6mm]}] (axis cs:0.1,0.15) -- (axis cs:0.12,0.02);
\end{axis}
\end{tikzpicture}
\caption{Сопротивление по формуле Блоха--Грюнайзена (расчёт): при
$T\gtrsim\theta_D/2$ рост почти линейный, при низких температурах~—
закон $T^5$. Примеси сдвигают кривую вверх на постоянную величину
$\rho_\text{ост}$ (правило Маттиссена)}
\label{fig:bloch-gr}
\end{figure}

\paragraph{Переходные и ферромагнитные металлы}
\begin{itemize}
\item У большинства переходных металлов (включая лантаноиды и актиноиды)
при $T<\theta_D$ сопротивление растёт как $T^n$ с $n=2$--$5{,}3$. Кроме
фононного рассеяния здесь работают другие механизмы: переходы
$s$-электронов на $d$-уровни (член $T^3$), рассеяние электронов на
электронах (член $T^2$); часто $\rho\approx AT^2+BT^5$.
\item Выше $\theta_D$ рост тоже нелинеен: у металлов IV и V групп
медленнее линейного, у металлов VI группы~— быстрее.
\item У ферромагнетиков (Fe, Co, Ni) добавочное рассеяние связано со
спонтанной намагниченностью; аномалии заметны примерно с $0{,}25\,\theta_K$
($\theta_K$~— точка Кюри), а выше точки Кюри сопротивление меняется как
у парамагнитных металлов (см. также п.~\ref{ssec:el-phase}).
\end{itemize}

\subsection{Эффект Холла}\label{ssec:hall}

Пусть по пластине шириной $b$ течёт ток плотности $j$, а
перпендикулярно пластине приложено магнитное поле $B$
(рис.~\ref{fig:hall}). Между боковыми гранями возникает разность
потенциалов~— \term{ЭДС Холла}. При не слишком сильных полях
\begin{keyeq}\label{eq:hall}
U_\text{Х}=R_\text{Х}\,B\,j\,b ,
\end{keyeq}
где $R_\text{Х}$~— \term{постоянная Холла} (м$^3$/Кл).

\begin{figure}[ht]
\centering
\begin{tikzpicture}[x={(1cm,0cm)}, y={(0.45cm,0.3cm)}, z={(0cm,1cm)}, scale=1.0, line join=round]
  \fill[fmBlue!10] (0,0,0) -- (5,0,0) -- (5,0,0.4) -- (0,0,0.4) -- cycle;
  \fill[fmBlue!18] (0,0,0.4) -- (5,0,0.4) -- (5,2.4,0.4) -- (0,2.4,0.4) -- cycle;
  \fill[fmBlue!6] (5,0,0) -- (5,2.4,0) -- (5,2.4,0.4) -- (5,0,0.4) -- cycle;
  \draw[atomLine] (0,0,0) -- (5,0,0) -- (5,0,0.4) -- (0,0,0.4) -- cycle
    (0,0,0.4) -- (0,2.4,0.4) -- (5,2.4,0.4) -- (5,0,0.4) (5,0,0) -- (5,2.4,0) -- (5,2.4,0.4);
  \fill[fmRed!60] (1.5,2.4,0.4) -- (3.5,2.4,0.4) -- (3.5,2.25,0.4) -- (1.5,2.25,0.4) -- cycle;
  \fill[fmBlue!60] (1.5,0.0,0.4) -- (3.5,0.0,0.4) -- (3.5,0.15,0.4) -- (1.5,0.15,0.4) -- cycle;
  \node[font=\scriptsize, fmRed] at (2.5,2.85,0.4) {$+\ +\ +$};
  \node[font=\scriptsize, fmBlue] at (2.5,-0.45,0.4) {$-\ -\ -$};
  \draw[-{Stealth[length=2.2mm]}, line width=1pt] (-1.2,1.2,0.2) -- (-0.1,1.2,0.2) node[pos=0.1, above, font=\small]{$\vb j$};
  \draw[-{Stealth[length=2.2mm]}, line width=1pt, fmGreen!60!black] (2.5,1.2,0.6) -- (2.5,1.2,1.9) node[above, font=\small]{$\vb B$};
  \draw[-{Stealth[length=2.2mm]}, fmRed, thick] (4.2,2.1,0.42) -- (4.2,0.4,0.42) node[right, font=\small]{$\vb E_\text{Х}$};
  \draw[dimline] (5.4,0,0) -- node[right, font=\small]{$b$} (5.4,2.4,0);
  \node[font=\small] at (0.4,2.9,0.4) {C};
  \node[font=\small] at (0.4,-0.5,0.4) {D};
\end{tikzpicture}
\caption{Эффект Холла (схема для электронов): сила Лоренца отклоняет
движущиеся носители к одной из граней, и возникает поперечное поле
$\vb E_\text{Х}$, уравновешивающее эту силу}
\label{fig:hall}
\end{figure}

\paragraph{Откуда берётся}
На движущийся электрон действует сила Лоренца $\vb F=-e[\vb v\vb B]$,
перпендикулярная скорости. Электроны отклоняются к одной грани, которая
заряжается отрицательно; противоположная, обеднённая электронами,~—
положительно. Возникает поперечное поле $E_\text{Х}$, и накопление
зарядов прекращается, когда сила поля уравновесит силу Лоренца:
$eE_\text{Х}=evB$. С $v=j/(ne)$ получаем $U_\text{Х}=bE_\text{Х}=bjB/(ne)$,
то есть
\begin{keyeq}\label{eq:RH}
R_\text{Х}=\frac1{ne}.
\end{keyeq}
Знак $R_\text{Х}$ совпадает со знаком носителей заряда: по эффекту Холла
находят и концентрацию, и знак носителей.

\paragraph{Опыт и зонная теория}
Сравним опыт с расчётом в предположении, что $n$ равна числу валентных
электронов (табл.~\ref{tab:hall}).

\begin{table}[ht]
\centering
\caption{Постоянная Холла металлов, $10^{-11}$~м$^3$/Кл}
\label{tab:hall}
\small
\begin{tabular}{@{}lccccccccc@{}}
\toprule
 & Na & Cs & Cu & Ag & Au & Bi & Be & Cd & Zn \\
\midrule
опыт & $-25$ & $-78$ & $-5{,}5$ & $-8{,}9$ & $-7{,}1$ & $-5\cdot10^4$ & $+24$ & $+6$ & $+10$ \\
расчёт $1/ne$ & $-24$ & $-73$ & $-7{,}4$ & $-10$ & $-10$ & $-4{,}1$ & $-2{,}5$ & $-6{,}4$ & $-4{,}6$ \\
\bottomrule
\end{tabular}
\end{table}

\begin{itemize}
\item Для щелочных и благородных металлов (I группа) согласие хорошее.
\item У Be, Cd, Zn знак \emph{положительный}~— как если бы ток переносили
положительные заряды. Объясняет это зонная теория. Если зона заполнена
меньше чем наполовину, электроны ведут себя нормально~— как частицы с
положительной эффективной массой и отрицательным зарядом, $R_\text{Х}<0$.
Если же зона заполнена почти целиком, удобнее описывать немногие
пустые состояния~— \emph{дырки}, которые ведут себя как частицы с
положительными массой и зарядом; тогда $R_\text{Х}>0$.
\item У Bi, Sb, As постоянная Холла аномально велика (у висмута почти в
$10^4$ раз больше, чем у меди): носителей вне почти заполненных зон
очень мало.
\end{itemize}

Эффект Холла~— один из основных методов изучения носителей заряда в
металлах и полупроводниках. На нём основаны датчики магнитного поля и
приборы для измерения токов; по тому же принципу (сила Лоренца на
проводящую среду с током) работают магнитные насосы для жидких металлов
(например, натриевого теплоносителя) и МГД-генераторы.
